\chapter{Énergie libre holographique-modulaire et mémoire informationnelle de la géométrie}
\label{chap:energia-libre-holografica}

\section{État bicouche orienté \(T_{A,y}\), normalisation \(\rho_y\) et séparation entre intensité et orientation de feuillet}

La réponse angulaire construite dans les chapitres précédents contient deux
données spectrales de nature complémentaire. La coordonnée
\(A_{\rm rad}\) fixe l'intensité commune aux deux feuillets ; la coordonnée
orientée
\[
 y=C^{*}_{\rm rad}
\]
fixe leur séparation. Soient \(P_+\) et \(P_-\) des projections orthogonales de
rang un sur un espace de Hilbert bidimensionnel \(\mathcal H_{\rm bi}\),
et définissons
\begin{equation}
 I_{\rm bi}=P_++P_-,
 \qquad
 R=P_+-P_-,
 \qquad
 R^2=I_{\rm bi}.
 \label{eq:rm-info-graduacion}
\end{equation}
L'opérateur angulaire positif est
\begin{equation}
 T_{A,y}
 =\exp(-A_{\rm rad}I_{\rm bi}-yR)
 =q_+P_++q_-P_-,
 \qquad
 q_\pm=\exp(-A_{\rm rad}\mp y).
 \label{eq:rm-info-operador-angular}
\end{equation}
Sa normalisation élimine exactement l'intensité scalaire et conserve l'
orientation relative :
\begin{equation}
 \boxed{
 \rho_y
 =\frac{T_{A,y}}{\Tr T_{A,y}}
 =\frac{\exp(-yR)}{2\cosh y}
 =p_+(y)P_++p_-(y)P_- ,}
 \label{eq:rm-info-rho-y}
\end{equation}
où
\begin{equation}
 p_+(y)=\frac{\ee^{-y}}{2\cosh y},
 \qquad
 p_-(y)=\frac{\ee^{y}}{2\cosh y}.
 \label{eq:rm-info-probabilidades-hoja}
\end{equation}
L'intensité \(A_{\rm rad}\) demeure dans \(\Tr T_{A,y}\) et dans le
déterminant de l'opérateur non normalisé ; l'état \(\rho_y\) représente le
secteur informationnel pur du choix de feuillet. Cette séparation type deux
lecteurs du même antécédent et évite d'attribuer à une probabilité
normalisée l'échelle constitutive retirée par la normalisation.

\begin{definition}[Hamiltonien holographique-modulaire bicouche]
L'hamiltonien holographique-modulaire fini de l'état orienté est
\begin{equation}
 \boxed{
 K_y=-\log\rho_y
 =\log(2\cosh y)I_{\rm bi}+yR.}
 \label{eq:rm-info-hamiltoniano-bicapa}
\end{equation}
Le terme scalaire normalise l'état et le terme \(yR\) conserve la
provenance de feuillet.
\end{definition}

La dénomination \emph{holographique-modulaire} déclare une fonction précise :
l'observable total publie l'occupation, tandis que le générateur modulaire
retient la distribution entre secteurs qui a rendu cette occupation possible. En
dimension finie, \(K_y\) est une matrice autoadjointe bien définie. Son
extension au système infini sera formulée au moyen d'un réseau local et de la
représentation GNS ; les deux résidences restent distinctes.

\section{Thermodynamique spectrale exacte des \(2\) feuillets}

La fonction de partition réduite et le moment orienté sont
\begin{equation}
 Z_{\rm bi}(y)=2\cosh y,
 \qquad
 \langle R\rangle_y=\Tr(\rho_yR)=-\tanh y.
 \label{eq:rm-info-particion-momento}
\end{equation}
L'entropie de von Neumann se calcule sans approximation :
\begin{equation}
 \boxed{
 S(\rho_y)
 =-\Tr(\rho_y\log\rho_y)
 =\log(2\cosh y)-y\tanh y.}
 \label{eq:rm-info-entropia-bicapa}
\end{equation}
Sa parité et sa variation satisfont
\begin{equation}
 S(\rho_{-y})=S(\rho_y),
 \qquad
 \frac{\dd}{\dd y}S(\rho_y)=-y\,\operatorname{sech}^2y.
 \label{eq:rm-info-paridad-entropia}
\end{equation}
Le déficit par rapport à l'état bicouche uniforme est déterminé par
\begin{equation}
 I_{\rm or}^{(2)}(y)
 =\log2-S(\rho_y)
 =y\tanh y-\log\cosh y
 \geq0.
 \label{eq:rm-info-deficit-orientado}
\end{equation}
L'égalité n'a lieu qu'en \(y=0\). Le déficit mesure l'information
retenue par la polarisation des feuillets, tandis que l'entropie mesure la
distribution probabiliste publiée.

\begin{theorem}[Entropie relative de deux orientations]
\label{thm:rm-info-relative-general}
Pour \(y,z\in\mathbb R\), les états fidèles de
\cref{eq:rm-info-rho-y} satisfont
\begin{equation}
 \boxed{
 D(\rho_y\Vert\rho_z)
 =\log\frac{\cosh z}{\cosh y}
 +(y-z)\tanh y.}
 \label{eq:rm-info-relative-general}
\end{equation}
En particulier, l'inversion de feuillet possède le coût exact
\begin{equation}
 \boxed{
 D(\rho_y\Vert\rho_{-y})=2y\tanh y.}
 \label{eq:rm-info-relative-inversion}
\end{equation}
\end{theorem}

\begin{proof}
De \(\log\rho_y=-\log(2\cosh y)I_{\rm bi}-yR\), il découle
\[
 \log\rho_y-\log\rho_z
 =\log\frac{\cosh z}{\cosh y}I_{\rm bi}+(z-y)R.
\]
En prenant l'espérance dans \(\rho_y\) et en utilisant
\(\langle R\rangle_y=-\tanh y\), on obtient
\cref{eq:rm-info-relative-general}. Substituer \(z=-y\) élimine le quotient
des cosinus hyperboliques et produit
\cref{eq:rm-info-relative-inversion}. L'expression est positive parce que
\(y\) et \(\tanh y\) ont le même signe.
\end{proof}

Le coût symétrisé de Jeffreys est
\begin{equation}
 D(\rho_y\Vert\rho_{-y})
 +D(\rho_{-y}\Vert\rho_y)
 =4y\tanh y.
 \label{eq:rm-info-jeffreys}
\end{equation}
La courbure locale de la famille exponentielle est donnée par
\begin{equation}
 \frac{\dd^2}{\dd y^2}\log Z_{\rm bi}(y)
 =\operatorname{Var}_{\rho_y}(R)
 =\operatorname{sech}^2y.
 \label{eq:rm-info-fisher}
\end{equation}
Ainsi, la même coordonnée orientée détermine trois quantités de parités
différentes : \(\langle R\rangle_y\) est impair ; \(S(\rho_y)\) est pair ; et la
distance entre orientations est positive et paire en tant que fonction du module
de \(y\). La géométrie, l'information et l'inversion conservent
des composantes distinctes d'un même état.

\begin{corollary}[Séparation de l'intensité et de l'orientation]
Le déterminant
\(\det T_{A,y}=\exp(-2A_{\rm rad})\) dépend de l'intensité angulaire,
tandis que le quotient \(q_-/q_+=\exp(2y)\), le moment
\(\langle R\rangle_y\) et
\(D(\rho_y\Vert\rho_{-y})\) dépendent de l'orientation. Par conséquent, le
lecteur pair qui conduit à \(D_A\) et le lecteur impair qui conduit à la mémoire
de feuillet sont conjointement reconstructibles et restent typés.
\end{corollary}

\section{Purification thermochamp double et conjugaison des feuillets}

Soit \((|+\rangle,|-\rangle)\) une base adaptée à
\((P_+,P_-)\), et soit \(\overline{\mathcal H}_{\rm bi}\) une copie
déclarée. Nous définissons
\begin{equation}
 |\operatorname{TFD}_y\rangle
 =\sqrt{p_+(y)}\,|+\rangle|\bar +\rangle
 +\sqrt{p_-(y)}\,|-\rangle|\bar -\rangle.
 \label{eq:rm-info-tfd-bicapa}
\end{equation}

\begin{proposition}[Purification bicouche]
\label{prop:rm-info-tfd-bicapa}
Le vecteur de \cref{eq:rm-info-tfd-bicapa} est normalisé et satisfait
\begin{equation}
 \Tr_{\overline{\mathcal H}_{\rm bi}}
 |\operatorname{TFD}_y\rangle
 \langle\operatorname{TFD}_y|
 =\rho_y.
 \label{eq:rm-info-tfd-partial}
\end{equation}
La permutation simultanée des étiquettes \(+\leftrightarrow-\)
transporte \(|\operatorname{TFD}_y\rangle\) en
\(|\operatorname{TFD}_{-y}\rangle\).
\end{proposition}

\begin{proof}
La somme des carrés des coefficients est
\(p_+(y)+p_-(y)=1\). La trace partielle élimine les termes croisés et
laisse \(p_+(y)P_++p_-(y)P_-\). L'échange des étiquettes permute
les deux coefficients, car \(p_+(-y)=p_-(y)\).
\end{proof}

La purification démontre une égalité des entropies au sein du système
doublé déclaré :
\begin{equation}
 S(\rho_y)
 =S\!\left(
 \Tr_{\mathcal H_{\rm bi}}
 |\operatorname{TFD}_y\rangle
 \langle\operatorname{TFD}_y|
 \right).
 \label{eq:rm-info-tfd-entropy}
\end{equation}
La copie thermochamp double possède ici une fonction mathématique : elle exhibe un
état pur dont la réduction produit exactement l'état bicouche. Son
identification à deux régions physiques d'un espace-temps requiert
une factorisation des observables et une interface géométrique supplémentaire.

La TFD bicouche et la TFD qutrit du
\cref{thm:modular-tfd} sont des constructions liées et distinctes. La
première purifie l'orientation de deux feuillets ; la seconde purifie les
trente-six liaisons qutrit des canaux \(90/120\). Leur confluence
réside dans l'antécédent angulaire et dans l'étiquette de
provenance, non dans une identification de leurs espaces de Hilbert.

\section{Géométrie totale et provenance modulaire}

Les opérateurs nombre \(N_{90}\) et \(N_{120}\) commutent. La géométrie
d'aire publie la somme des occupations,
\begin{equation}
 \widehat{\mathcal A}
 =\gammaH a_0\ell_P^2(N_{90}+N_{120}),
 \label{eq:rm-info-area-total}
\end{equation}
tandis que l'hamiltonien modulaire de bord conserve leur provenance,
\begin{equation}
 K_A-cI
 =\Dmod(90N_{90}+120N_{120}).
 \label{eq:rm-info-hhm-provenance}
\end{equation}
Définissons les observables adimensionnels
\begin{equation}
 N=N_{90}+N_{120},
 \qquad
 D=90N_{90}+120N_{120}.
 \label{eq:rm-info-ND}
\end{equation}

\begin{theorem}[Reconstruction conjointe des canaux]
\label{thm:rm-info-reconstruction-90-120}
La paire \((N,D)\), et donc la paire
\((\widehat{\mathcal A},K_A)\) une fois ses échelles fixées, reconstruit
exactement les deux opérateurs d'occupation :
\begin{equation}
 \boxed{
 N_{90}=4N-\frac{D}{30},
 \qquad
 N_{120}=\frac{D}{30}-3N.}
 \label{eq:rm-info-inverse-90-120}
\end{equation}
\end{theorem}

\begin{proof}
Le système linéaire a pour matrice
\(\left(\begin{smallmatrix}1&1\\90&120\end{smallmatrix}\right)\)
de déterminant \(30\). Son inverse donne
\cref{eq:rm-info-inverse-90-120}. La commutativité des opérateurs
nombre permet d'appliquer l'identité dans leur calcul fonctionnel conjoint.
\end{proof}

Le résultat formule une holographie de reconstruction concrète. L'
aire détermine la quantité géométrique totale ; l'hamiltonien modulaire
détermine la provenance pondérée ; la paire retrouve la distribution
sectorielle. Deux états de même aire peuvent porter des histoires
distinctes, et cette différence reste visible dans \(K_A\). Le coefficient
\(30\) qui rend le système inversible est le même défaut de complétion
entre les canaux \(90\) et \(120\).

La reconstruction exige de conserver les deux observables. Une compression qui
ne retiendrait que \(N\) identifierait toutes les paires
\((N_{90},N_{120})\) de somme constante. Une compression qui ne retiendrait que
\(D\) identifierait des paires de même poids modulaire. La paire conjointe possède un
rang de deux et conserve la généalogie sectorielle pertinente.

\section{1ʳᵉ loi locale et bilan fini d'entropie relative}

La formulation finie emploie l'entropie relative des états opératoriels
au sens d'Araki \cite{Araki1976}.

Soit \(\rho_A\) l'état qutrit de bord de
\cref{eq:modular-product-state} et soit
\(K_A=-\log\rho_A\). La décomposition d'aire démontrée dans le
\cref{thm:modular-first-law} prend la forme
\begin{equation}
 K_A-cI
 =\beta_{90}\mathcal A_{90}
 +\beta_{120}\mathcal A_{120},
 \qquad
 \frac{\beta_{120}}{\beta_{90}}=\frac43.
 \label{eq:rm-info-K-area}
\end{equation}
Pour une courbe différentiable d'états normalisés passant par
\(\rho_A\), la variation tangente satisfait
\begin{equation}
 \boxed{
 \delta S
 =\beta_{90}\,\delta\langle\mathcal A_{90}\rangle
 +\beta_{120}\,\delta\langle\mathcal A_{120}\rangle.}
 \label{eq:rm-info-first-law}
\end{equation}
Le quotient \(4/3\) enregistre les poids informationnels différents des deux
secteurs ; un même accroissement d'aire totale peut produire des variations
entropiques différentes selon son canal de provenance.

La formulation finie conserve le terme que la linéarisation élimine.
Si \(\sigma\) est un état dont le support est contenu dans celui de
\(\rho_A\), alors
\begin{equation}
 \boxed{
 \Delta S
 =\beta_{90}\Delta\langle\mathcal A_{90}\rangle
 +\beta_{120}\Delta\langle\mathcal A_{120}\rangle
 -D(\sigma\Vert\rho_A).}
 \label{eq:rm-info-finite-law}
\end{equation}

\begin{proof}
La définition de l'entropie relative fournit
\[
 D(\sigma\Vert\rho_A)
 =-S(\sigma)+\Tr(\sigma K_A).
\]
Pour \(\rho_A\), la même expression vaut zéro. En soustrayant les deux
identités, on obtient
\(\Delta S=\Delta\langle K_A\rangle-D(\sigma\Vert\rho_A)\).
Le terme scalaire de \(K_A\) s'annule entre états normalisés, et
\cref{eq:rm-info-K-area} conclut la preuve.
\end{proof}

La loi locale exprime la dérivée de l'entropie dans l'état de
référence. La loi finie exprime le bilan exact entre réponse
géométrique, changement entropique et distance informationnelle. La positivité de
\(D\) quantifie la convexité qui sépare le déplacement fini de son
approximation tangente.

\section{Énergie libre holographique-modulaire d'horizon}

Considérons une réalisation d'horizon dans laquelle un rayon d'aire \(R\)
détermine
\begin{equation}
 E_\Lambda(R)=\frac{c^4R}{2G},
 \qquad
 k_BT_H(R)=\frac{\hbar c}{2\pi R},
 \qquad
 \frac{S_H(R)}{k_B}=\frac{\pi R^2}{\ell_P^2}.
 \label{eq:rm-info-horizon-data}
\end{equation}
L'identité \(\ell_P^2=G\hbar/c^3\) donne
\begin{equation}
 \boxed{E_\Lambda(R)=T_H(R)S_H(R).}
 \label{eq:rm-info-horizon-saturation}
\end{equation}
Pour une perturbation finie \(\sigma\) de l'état de référence
\(\rho_A\), écrivons
\begin{equation}
 \Delta k_A=\Tr((\sigma-\rho_A)K_A),
 \qquad
 \Delta s_A=S(\sigma)-S(\rho_A),
 \label{eq:rm-info-delta-k-s}
\end{equation}
et définissons les grandeurs de réponse
\begin{equation}
 \mathcal E_R(\sigma)
 =E_\Lambda(R)+k_BT_H(R)\Delta k_A,
 \qquad
 \mathcal S_R(\sigma)
 =S_H(R)+k_B\Delta s_A.
 \label{eq:rm-info-response-energy-entropy}
\end{equation}

\begin{theorem}[Identité d'énergie libre holographique-modulaire]
\label{thm:rm-info-free-energy}
Sous la réalisation d'horizon
\cref{eq:rm-info-horizon-data}, pour tout \(\sigma\) dont le support est contenu
dans celui de \(\rho_A\),
\begin{equation}
 \boxed{
 \mathcal E_R(\sigma)-T_H(R)\mathcal S_R(\sigma)
 =k_BT_H(R)D(\sigma\Vert\rho_A)
 \geq0.}
 \label{eq:rm-info-free-energy}
\end{equation}
\end{theorem}

\begin{proof}
La saturation de référence
\cref{eq:rm-info-horizon-saturation} annule le terme
\(E_\Lambda-T_HS_H\). Par définition,
\[
 \mathcal E_R-T_H\mathcal S_R
 =k_BT_H(\Delta k_A-\Delta s_A).
\]
L'identité finie d'entropie relative donne
\(\Delta k_A-\Delta s_A=D(\sigma\Vert\rho_A)\), et l'inégalité de Klein
donne la positivité.
\end{proof}

Dans la clôture \(R=\ell_P\), le quantum de Planck
\(E_P=\hbar c/\ell_P\) produit
\begin{equation}
 k_BT_H(\ell_P)=\frac{E_P}{2\pi},
 \qquad
 E_\Lambda(\ell_P)=\frac{E_P}{2},
 \label{eq:rm-info-planck-horizon-data}
\end{equation}
et donc
\begin{equation}
 \boxed{
 \mathcal E_{\ell_P}(\sigma)
 -T_H(\ell_P)\mathcal S_{\ell_P}(\sigma)
 =\frac{E_P}{2\pi}D(\sigma\Vert\rho_A).}
 \label{eq:rm-info-planck-free-energy}
\end{equation}
L'identité structurelle
\begin{equation}
 E_P=k_B\Theta_{\rm clk}\log3
 \label{eq:rm-info-planck-trit}
\end{equation}
transforme le coefficient du défaut en
\begin{equation}
 \frac{E_P}{2\pi}
 =\frac{k_B\Theta_{\rm clk}\log3}{2\pi}.
 \label{eq:rm-info-trit-free-energy}
\end{equation}
L'horizon de référence sature le bilan ; l'écart fini est
mesuré par le coût informationnel de l'abandon de sa généalogie modulaire.

La réalisation satisfait aussi
\begin{equation}
 E_\Lambda(\ell_P)t_P=\frac{\hbar}{2},
 \qquad
 E_\Lambda(\ell_P)T_{108}=\frac h2,
 \label{eq:rm-info-half-action}
\end{equation}
lorsque \(t_P=\ell_P/c\) et \(T_{108}=2\pi t_P\). La demi-action
résulte du couplage entre périodicité thermique et géométrie d'aire.
Cette égalité reste séparée de l'identité d'énergie libre : la
première fixe le quantum de l'horizon de référence ; la seconde quantifie
les perturbations de l'état modulaire.

\section{Confluence entre quantum d'aire et quantum informationnel}

La transition hexade--octade publie le quantum géométrique
\begin{equation}
 \delta\mathcal A
 =\gammaH a_0\ell_P^2,
 \label{eq:rm-info-area-quantum}
\end{equation}
et l'horloge CT108 publie le quantum informationnel
\begin{equation}
 E_P=k_B\Theta_{\rm clk}\log3.
 \label{eq:rm-info-energy-quantum}
\end{equation}
Leur quotient définit la tension typée
\begin{equation}
 \boxed{
 \mathcal T_{I/A}
 =\frac{E_P}{\delta\mathcal A}
 =\frac{k_B\Theta_{\rm clk}\log3}
 {\gammaH a_0\ell_P^2}.}
 \label{eq:rm-info-tension-information-area}
\end{equation}
L'équation assigne une conversion entre une décision ternaire et une
unité de géométrie. Le numérateur appartient à la branche
action--horloge--information ; le dénominateur appartient à la branche
hexade--octade--aire. Toutes deux descendent du même état enrichi et
confluent dans la gravité après avoir conservé leurs domaines.

La paire
\begin{equation}
 \left(
 \widehat{\mathcal A},K_A
 \right)
 \label{eq:rm-info-area-K-pair}
\end{equation}
reconstruit la quantité et la provenance ; la paire
\begin{equation}
 \left(
 \delta\mathcal A,E_P
 \right)
 \label{eq:rm-info-area-energy-pair}
\end{equation}
relie la géométrie élémentaire et le coût informationnel. La première paire est
spectrale ; la seconde est dimensionnelle. Ensemble, elles décrivent comment une structure
discrète d'orientation est publiée comme aire, énergie et entropie.

\section{Réseau local UHF et générateur modulaire dans la représentation GNS}

Soit \((\mathcal A_m,\iota_m)\) un réseau d'algèbres matricielles locales
avec
\begin{equation}
 \mathcal A_m\simeq M_{9^m}(\mathbb C),
 \qquad
 \iota_m(a)=a\otimes I_9,
 \qquad
 \mathcal A_{\rm UHF}=\varinjlim_m\mathcal A_m.
 \label{eq:rm-info-uhf-net}
\end{equation}
Une famille d'états locaux \(\omega_m\) définit un état global
\(\omega\) lorsqu'elle satisfait
\begin{equation}
 \omega_{m+1}\circ\iota_m=\omega_m.
 \label{eq:rm-info-state-compatibility}
\end{equation}
À chaque coupure finie, on peut écrire
\begin{equation}
 \omega_m(a)=\Tr(\rho_m a),
 \qquad
 K_m=-\log\rho_m.
 \label{eq:rm-info-local-densities}
\end{equation}
La famille \((K_m)\) constitue la résidence locale de l'hamiltonien
holographique-modulaire.

Dans la représentation GNS
\((\pi_\omega,\mathcal H_\omega,\Omega_\omega)\), la théorie modulaire s'
exprime au moyen de l'opérateur modulaire \(\Delta_\omega\) et du groupe
\begin{equation}
 \sigma_t^\omega(x)
 =\Delta_\omega^{\,it}x\Delta_\omega^{-it},
 \qquad
 x\in\pi_\omega(\mathcal A_{\rm UHF})''.
 \label{eq:rm-info-gns-modular-flow}
\end{equation}
Le générateur \(-\log\Delta_\omega\) peut être non borné et appartient à la
représentation standard ; une matrice de densité globale de classe trace dans
l'UHF ne fait pas partie, en général, de cette construction.

L'espérance traciale nonadique
\begin{equation}
 E_m=\operatorname{id}\otimes\tau_9:
 \mathcal A_{m+1}\longrightarrow\mathcal A_m
 \label{eq:rm-info-tracial-expectation}
\end{equation}
préserve l'unité et constitue le lecteur canonique de compression
matricielle. La compatibilité d'un état modulaire exige en outre
\begin{equation}
 \omega_m\circ E_m=\omega_{m+1}
 \quad\hbox{dans le sous-espace déclaré,}
 \label{eq:rm-info-expectation-state}
\end{equation}
ou la condition équivalente adaptée à l'inclusion utilisée. La
tracialité du lecteur ne transforme pas automatiquement un produit de Gibbs non
tracial en une famille compatible.

Le \cref{thm:modular-typeIII} démontre, sous persistance fidèle des deux
poids \(90\Dmod\) et \(120\Dmod\), densité asymptotique positive des deux
canaux, factorialité et absence de quotients supplémentaires, que le groupe des
rapports est
\begin{equation}
 90\Dmod\mathbb Z+120\Dmod\mathbb Z
 =30\Dmod\mathbb Z
 \label{eq:rm-info-ratio-group}
\end{equation}
et que la fermeture ITPFI hyperfinie est de type
\begin{equation}
 \boxed{
 III_{\lambda_\partial},
 \qquad
 \lambda_\partial=\exp(-30\Dmod).}
 \label{eq:rm-info-type-iii}
\end{equation}
L'énoncé classifie le réseau produit spécifié. Une modification du
groupe des rapports, la disparition asymptotique d'un canal ou la perte de
factorialité oblige à recalculer le type. En particulier, la fermeture traciale
II\(_1\) de l'UHF et la fermeture modulaire de type III correspondent à des états et à des
représentations différents ; toutes deux sont des lecteurs opératoriels du système
nonadique et aucune ne remplace l'antécédent APP--TRIT--TPK.

\section{Identités exactes de la clôture gravitationnelle et conditions de la réalisation tensorielle}

L'architecture de ce chapitre distingue les résultats suivants :
\begin{enumerate}[label=\textup{(\roman*)}]
 \item La normalisation bicouche, l'hamiltonien
 \(K_y\), l'entropie
 \(S(\rho_y)\) et l'identité
 \(D(\rho_y\Vert\rho_{-y})=2y\tanh y\) sont une
 \status{nouvelle formalisation exacte} de la réponse angulaire déjà
 construite.
 \item La TFD bicouche est une \status{nouvelle formalisation exacte} dans l'
 espace doublé déclaré. L'interprétation de la copie comme une autre
 région physique requiert une factorisation géométrique ultérieure.
 \item La reconstruction \(90/120\), la TFD qutrit, la première loi locale,
 l'identité finie et la classification modulaire sous leurs hypothèses sont
 \status{resultados exactos} de la construction opératorielle précédente.
 \item L'égalité d'horizon
 \(E_\Lambda=T_HS_H\) est exacte dans la carte de Sitter déclarée. Sa
 composition avec l'état qutrit pour former
 \cref{eq:rm-info-free-energy} est une
 \status{interface physique conditionnée} par cette réalisation.
 \item L'entrelaceur entre microétats d'horizon, TFD de bord et
 sous-espaces propres hyperboliques est classifié par spectre et multiplicité
 dans le \cref{thm:ivv-horizon-tfd-perron-classification} : la coïncidence
 typée produit l'équivalence unitaire et son échec démontre l'obstruction.
\end{enumerate}

\begin{proofstatusbox}
Les identités spectrales de la famille \(\rho_y\), sa purification,
la reconstruction linéaire des secteurs \(90/120\) et le bilan fini
avec entropie relative sont exacts en dimension finie. La classification de
type III est exacte sous les quatre hypothèses formulées dans cette construction.
L'énergie libre gravitationnelle est affirmée dans la réalisation d'horizon de
\cref{eq:rm-info-horizon-data} ; elle conserve la résidence propre de l'état
modulaire et maintient la dynamique cosmologique comme une promotion ultérieure.
\end{proofstatusbox}

\begin{falsifierbox}
Réfutent les affirmations correspondantes : un spectre de \(\rho_y\) dont la somme ne
vaut pas un ; un calcul de
\(D(\rho_y\Vert\rho_{-y})\) différent de \(2y\tanh y\) ; une trace partielle de
la TFD qui ne produit pas \(\rho_y\) ; deux paires distinctes
\((N_{90},N_{120})\) ayant les mêmes \((N,D)\) ; une perturbation finie
qui viole \cref{eq:rm-info-finite-law} ; une carte d'horizon qui ne satisfait pas
\(E_\Lambda=T_HS_H\) ; ou un réseau infini dont le groupe des rapports diffère de
\(30\Dmod\mathbb Z\) sous les hypothèses déclarées. Réfute également une
promotion opératorielle le remplacement du réseau local par une matrice de densité
globale inexistante ou l'identification de la fermeture traciale II\(_1\) à la
fermeture modulaire de type III.
\end{falsifierbox}
